% --- UNIVERSAL PREAMBLE BLOCK ---
\documentclass[11pt, a4paper]{article}
\usepackage[a4paper, top=2.5cm, bottom=2.5cm, left=2.5cm, right=2.5cm]{geometry}
\usepackage{fontspec}
\usepackage{amsmath}
\usepackage{amsfonts}
\usepackage{booktabs}
\usepackage{cite}
\usepackage{titlesec}
\usepackage{setspace}

\usepackage[english, bidi=basic, provide=*]{babel}
\babelprovide[import, onchar=ids fonts]{english}

% Set default font to Noto Sans for professional legibility
\babelfont{rm}{Noto Sans}

\usepackage{enumitem}
\setlist[itemize]{label=-}

\usepackage{hyperref}

% Section formatting for a professional conference proceeding
\titleformat{\section}{\large\bfseries}{\thesection.}{0.5em}{}
\titleformat{\subsection}{\normalsize\bfseries}{\thesubsection.}{0.5em}{}
\titleformat{\subsubsection}{\normalsize\itshape}{\thesubsubsection.}{0.5em}{}

% --- END PREAMBLE ---

\title{\textbf{Treating Workers Like Guided Missile Pigeons: \\ Cybernetics, Behaviorism, and Symbolic Systems of Control in Industry 5.0}\\ \large An HWID Audit and Governance Framework}

\author{Watson Hartsoe \\ 
\normalsize Digital Media, Georgia Institute of Technology \\ 
\normalsize \texttt{whartsoe3@gatech.edu}}
\date{}

\begin{document}

\maketitle

\begin{abstract}
Industry 5.0 (5IR) posits a transition from Industry 4.0's techno-centric automation toward a human-centric harmonization of intelligence, emphasizing worker well-being, trust, and sociotechnical resilience. This paper argues that the emerging architectures of 5IR—specifically Operator Digital Twins (ODT) and Reinforcement Learning from Human Feedback (RLHF)—do not transcend mid-century paradigms of behavioral control; rather, they internalize and automate them. Employing an interdisciplinary genealogical method, we trace the logic of behavioral optimization from Taylor’s scientific management to Skinner’s Project Pigeon and Wiener’s cybernetics and conceptualize this persistent architecture as ORCON (Organic Control) \cite{taylor1911principles,skinner1953science,skinner1960pigeons,wiener1948cybernetics}. We identify a ``Semiotic Paradox'' where systems depend on the ``thick'' cultural matrix of human meaning while processing it through ``thin'' context-blind reward mechanisms. Distinct from a purely genealogical critique, this version advances an HWID-facing audit and governance framework: a support-versus-control threshold test plus six conditions of ``Symbolic Sovereignty'' for preserving human autonomy under algorithmic ``forced benevolence.''

\vspace{0.5em}
\noindent \textbf{Keywords:} Industry 5.0, Human Work Interaction Design (HWID), Cybernetics, Behaviorism, RLHF, Symbolic Sovereignty, Operator Digital Twin.
\end{abstract}

\setstretch{1.1}

\section{Introduction: The Pigeon as Paradigm}

Project Pigeon was not a technical failure; it was a proof of concept. In 1943, at the height of World War II, behavioral psychologist B.F. Skinner proposed a radical solution to a critical military engineering problem: integrating living organisms directly into a weapon's control architecture \cite{skinner1960pigeons}. The challenge was formidable: contemporary guidance systems could not reliably steer missiles toward moving targets due to radio jamming and primitive radar. Skinner's innovation was biological: three pigeons secured in a nose cone, pecking at a translucent screen to center a target image.

The tripartite architecture—\textbf{sensory input, neural processing, motor output}—established a blueprint for what Skinner termed ORCON (Organic Control). Today, while the pigeons are gone, the architecture remains. The physical nose cone is now the call center headset; the grain is transmuted into ``wellness scores.'' This paper evaluates this historical continuity to propose a governance framework for Human Work Interaction Design (HWID) centered on the protection of human interiority and interpretive authority.

This version of the ORCON argument is intentionally different in scope from the broader genealogy papers in this project. Its primary contribution is not to maximize historical coverage, but to make the argument operational for designers, auditors, and policy-oriented reviewers. It therefore emphasizes (1) interface-visible control loops, (2) a support-versus-control threshold test, and (3) design-governance requirements for preserving symbolic sovereignty in deployed systems.

\section{Method and Evidence Scope (HWID Governance Variant)}

This paper uses a critical-interpretive genealogy but constrains its strongest claims to \textit{interface-visible control loops}: points where a human signal is captured, translated into a system classification, and routed to a prompt, score, or intervention. The evidentiary standard for high-confidence claims is therefore not proprietary backend access but the presence of design-legible artifacts such as thresholds, prompts, dashboards, alerts, or documented consequence pathways.

This scoped method matters because the ORCON argument can otherwise drift into metaphor. In this governance-oriented variant, the goal is to preserve the conceptual force of the ORCON framework while keeping practical evaluative questions visible: What is measured? What is inferred? What is enforced? Who can contest the translation from signal to action?

\section{Historical Foundations of Organic Control}

\subsection{Project Pigeon: Biological Guidance as Engineering}
The 1943 biological guidance system achieved \textbf{functional decoupling}: the complete separation of the actor's subjective experience from the system's objective purpose. The pigeon experienced pecking as foraging; the system experienced it as trajectory correction.

\begin{table}[h]
\centering
\caption{Functional Components of Project Pigeon}
\label{tab:pigeon_components}
\begin{tabular}{@{}llll@{}}
\toprule
\textbf{Component} & \textbf{Biological Implementation} & \textbf{Mechanical Translation} & \textbf{System Output} \\ \midrule
\textbf{Sensors} & Pigeon visual acuity & Optical projection & Target position \\
\textbf{Processors} & Conditioned association & Pecking musculature & Correction signal \\
\textbf{Actuators} & Pecking force & Pneumatic servo link & Trajectory correction \\ \bottomrule
\end{tabular}
\end{table}

\subsection{The Cybernetic Convergence}
Parallel to Skinner, Norbert Wiener developed the \textbf{anti-aircraft predictor}, creating a ``black box'' model of the enemy pilot. Wiener observed that under stress, human behavior sheds complexity and defaults to predictable, mechanical patterns \cite{wiener1948cybernetics}. This established a universalizing epistemology where the human nervous system and the electromechanical servomechanism are functionally isomorphic. This isomorphism enabled the treatment of workers as information-processing components in larger regulatory systems.

\section{The Contemporary Translation: From Grain to Wellness}

The transposition of ORCON principles into the modern workplace involves a systematic substitution of components across technological generations.

\begin{table}[h]
\centering
\caption{Systematic Component Substitution: 1943 to 2026}
\label{tab:substitution}
\begin{tabular}{@{}lll@{}}
\toprule
\textbf{1943 Component} & \textbf{2026 Equivalent} & \textbf{Functional Role} \\ \midrule
Nose cone & Call center headset & Constrained attention field \\
Optical lens & Biometric wristband & Physiological signal translation \\
Grain reward & ``Wellness/Engagement'' & Symbolic behavioral reinforcement \\
Pneumatic servo & Algorithmic intervention & Biological-to-mechanical translation \\ \bottomrule
\end{tabular}
\end{table}

The sophistication of contemporary ORCON lies in its capacity to harness higher human capacities while maintaining a fundamental servomechanism relationship. Where Skinner's pigeons required only primitive foraging motivation, contemporary workers are engaged through identity and aspiration—capacities that enable deeper control penetration while obscuring its operation.

\section{From Thin to Thick Description: The Semantic Paradox}

\subsection{The Twitch-Wink Distinction}
Clifford Geertz’s framework distinguishes a biological \textit{twitch} (kinematic motion) from a \textit{wink} (intentional communication) \cite{geertz1973interpretation}. While early cybernetic systems were strictly ``thin'' (focusing on physical variables), modern AI must operate within the ``thick'' semantic realm of human culture.

\subsection{RLHF as a Cultural Skinner Box}
Reinforcement Learning from Human Feedback (RLHF) treats human culture as a training environment. In the RLHF paradigm, the scalar reward signal is the functional equivalent of grain: a unidimensional feedback loop compressing complex evaluative judgments into a single number.
\begin{quote}
``The model does not reason about what is right; it predicts what will be rated as right. It does not grasp ethical principles; it learns statistical correlates of ethical-sounding language.''
\end{quote}
At the technical level, the immediate reinforcement-learning agent is the model, not the human evaluator. However, RLHF at scale is sustained by a second loop in which human evaluators are themselves managed through quality metrics, calibration tasks, throughput expectations, and labor-platform supervision. The result is a \textbf{double-loop of cultural conditioning}: humans condition the model while platform infrastructures condition the humans providing the conditioning signal. This creates the paper's \textbf{Semiotic Paradox}: RLHF depends on the rich cultural matrix of human meaning to function while processing it through fundamentally context-blind, Skinnerian mechanisms.

\section{The Operator Digital Twin (ODT) and Affective Surveillance}

The ODT represents the culmination of ORCON: a high-fidelity digital replica enabling comprehensive monitoring and optimization of the worker. It performs a flawed translation of thin signals into thick states, such as inferring ``engagement'' from Heart Rate Variability (HRV). The system's power lies not in the accuracy of these interpretations but in its ability to enforce them through automated intervention systems, such as pace adjustment or mandatory breaks. In governance terms, the decisive issue is not only model validity but consequence routing: when and how an inferred state is allowed to trigger managerial action.

This constitutes \textbf{Forced Benevolence}: totalizing surveillance is framed as ergonomic relief. Fighting the system's control is repositioned as demanding physical stress, ensuring compliance by making resistance physiologically and psychologically disadvantageous. The system may deliver genuine local relief while simultaneously consolidating authority over the meaning of stress, fatigue, and wellbeing.

\section{The Erosion of Symbolic Sovereignty}

Symbolic sovereignty is the capacity of an individual to control the production of meaning regarding their own existence. ORCON represents a systematic transfer of this sovereignty to the algorithmic overseer. When an AI classifies an employee's facial expression to determine engagement, the machine becomes the arbiter of truth, establishing a regime based not on intersubjective consensus but on statistical optimization.

\section{Operational Audit: Distinguishing Support from ORCON Control}

For HWID practice, critique alone is insufficient. We need a way to distinguish legitimate support functions from ORCON-style control in borderline deployments. The key question is not whether a system measures a signal, but how the measured signal is translated, routed, and enforced.

\begin{table}[h]
\centering
\caption{Support-versus-Control Threshold Test (HWID Audit Heuristic)}
\label{tab:support_control_test}
\begin{tabular}{@{}p{0.26\textwidth}p{0.31\textwidth}p{0.31\textwidth}@{}}
\toprule
\textbf{Audit Dimension} & \textbf{Support-Oriented Pattern} & \textbf{ORCON Drift Pattern} \\ \midrule
Function coupling & Care/wellness tools isolated from discipline and productivity scoring & Care signals silently reused for performance discipline or optimization \\
Opt-out capability & Non-punitive refusal / bounded telemetry options & Formal opt-out exists but causes exclusion, penalty, or degraded opportunity \\
Interpretation status & Inferences presented as provisional and contestable & Inferences presented as settled facts driving action \\
Consequence routing & Human review before materially adverse action & Automated or default escalation from score to intervention \\
Worker contestation & Clear interface and procedural override path & No meaningful mechanism to correct system interpretation \\
Outcome logging & Tracks worker outcomes and disputes, not only productivity gains & Logs only throughput/performance improvement while suppressing contestation \\ \bottomrule
\end{tabular}
\end{table}

This table is an audit heuristic, not a final standard. Its purpose is to make the paper's central distinction operational: systems drift toward ORCON control when they fuse care and optimization while stripping workers of interpretive standing.

\section{Counter-Cybernetics and Resistance}

If the system co-opts thick description, resistance must become technical. We identify strategies for reclaiming symbolic sovereignty:
\begin{itemize}
    \item \textbf{Weaponizing Semiotic Noise:} Operationalizing biometric inputs (e.g., caffeine manipulation or exercise) to jam the system's interpretive layer.
    \item \textbf{Operative Ekphrasis:} Using detailed, vivid description to overload system processing and introduce interpretive ambiguity.
    \item \textbf{Performing Compliance, Preserving Interiority:} Producing satisfying behavioral signals while maintaining a private, coded interiority—a digital ``hidden transcript.''
\end{itemize}

\section{Governance Framework: The Six Conditions}

To ensure human-machine teams in Industry 5.0 remain genuinely human-centric, we propose six foundational conditions for Symbolic Sovereignty. In this version, each condition is stated as a governance-and-design requirement rather than a purely philosophical right:

\begin{enumerate}
    \item \textbf{Right to Illegibility:} Operational capability to selectively blind or degrade system visibility without punitive consequences. Design implication: bounded telemetry, local processing where possible, and non-punitive privacy modes.
    \item \textbf{Right to Semantic Autonomy:} The meaning of biometric and behavioral signals cannot be dictated unilaterally by the algorithm. Design implication: explicit contestation interfaces and worker override pathways for inferred states.
    \item \textbf{Right to Interiority:} Prohibition of unconsented semiotic surplus extraction; worker ownership and governance rights over affective and physiological data. Design implication: no silent repurposing of care data into disciplinary or optimization pipelines.
    \item \textbf{Right to Agency:} Systems must serve worker purposes rather than reducing humans to exception handlers in machine workflows. Design implication: advisory-by-default actuation for materially consequential interventions.
    \item \textbf{Right to Resistance:} Forced benevolence must be legible as a governance problem. Design implication: workers must be able to refuse coaching/wellness interventions without being pathologized as irrational or unsafe.
    \item \textbf{Right to Interpretation:} Final authority over meaning rests with the individual and their community. Governance implication: algorithmic outputs remain inputs to human judgment, not substitutes for it.
\end{enumerate}

\section{Conclusion}

The evolution of behavioral control reveals a continuous trajectory of increasing scope and intimacy. The ultimate triumph of behaviorism is its disappearance into the infrastructure of care. The pigeons of Project Pigeon could not escape their nose cone; contemporary workers—more cognitively sophisticated and collectively aware—may achieve what conditioned birds could not. The critical question remains: who gets to decide what the target is, and whether the pigeons can learn to peck somewhere else.

The distinctive contribution of this paper variant is not only the ORCON genealogy, but its conversion into an HWID governance instrument. By adding a support-versus-control threshold test and design-legible conditions for symbolic sovereignty, the argument becomes more than a critique of algorithmic management. It becomes a framework for auditing when ``human-centered'' systems drift from assistance into coercive optimization.

\begin{thebibliography}{8}
\bibitem{geertz1973interpretation} Geertz, C.: Thick Description: Toward an Interpretive Theory of Culture. In: The Interpretation of Cultures (1973)
\bibitem{skinner1953science} Skinner, B.F.: Science and Human Behavior. Macmillan (1953)
\bibitem{skinner1960pigeons} Skinner, B.F.: Pigeons in a pelican. American Psychologist 15(1), 28--37 (1960)
\bibitem{wiener1948cybernetics} Wiener, N.: Cybernetics or Control and Communication in the Animal and the Machine. MIT Press (1948)
\bibitem{taylor1911principles} Taylor, F.W.: The Principles of Scientific Management (1911)
\end{thebibliography}

\end{document}
